% slides_zh_frames.tex -- the Chinese screen of every slide, frame for frame
% with slides.tex. tools/build_zh_deck.py joins each block with the
% matching \note{} from slides.tex to make slides_zh.tex, so the notes are
% written once. Terms stay English, with the Chinese in brackets.
%%% FRAME 1 {}
\vspace*{4mm}
{\usebeamerfont{title}\usebeamercolor[fg]{frametitle}Green Machines（植物）+ Metallic materials（金属）}\par
\vspace{1mm}{\color{rule}\rule{\textwidth}{0.8pt}}\par\vspace{2mm}
Year 9 科学 \textperiodcentered{} 考前冲刺 \textperiodcentered{} 2 小时\par\vspace{5mm}
\begin{tabular}{@{}l@{\hspace{6mm}}l@{}}
0:00 & 第一课 \textperiodcentered{} Green Machines（植物）\\
0:50 & 第二课 \textperiodcentered{} Metallic materials（金属）\\
1:40 & Mini mock 小模考（两个单元）\\
\end{tabular}
\answers{%
本页没有题目。0:00 开始：先发 Classwork；Mini mock 到 1:40 再发。}
\fref{中文版}
%%% FRAME 2 {}
\lessondivider{第一课 \textperiodcentered{} Green Machines（植物）}{学习目标：用考试用的英文词，解释植物怎样制造食物、运输水分、产生新植物。}
\answers{%
本页没有题目。这一小时：Loop 1 繁殖 → Loop 2 光合作用 → Loop 3 运输。}
\fref{中文版}
%%% FRAME 3 {Do Now 课前小测}
\qline{1} 花的哪一部分产生 pollen（花粉）？
\qline{2} 花粉的 nucleus（细胞核）和 ovule（胚珠）的细胞核融合，这叫 pollination（授粉）还是 fertilisation（受精）？
\qline{3} 种子 germinate（萌发）需要哪三个条件？
\qline{4} 叶子吸收哪种气体来制造 glucose（葡萄糖）？
\qline{5} 哪种管道把水沿着茎往上运？
\answers{%
1  anther（花药）\par
2  fertilisation（受精）\par
3  water, warmth, oxygen（三个都要说）\par
4  carbon dioxide（二氧化碳）\par
5  xylem（木质部）}
\fref{Green Machines p.2}
%%% FRAME 4 {The one idea 一个核心思想}
\vspace{6mm}
\FigGreenIdea\par\vspace{8mm}
每道 Green Machines 的题，都在问这三件事中的一件。\par\vspace{2mm}
先判断是哪一件，再用那一件的英文词作答。
\answers{%
问：用碘液测 starch 属于哪件事？答：MAKE food，因为 starch 是储存的 glucose。}
\fref{Green Machines p.2}
%%% FRAME 5 {Quick review 快速回忆：花的结构}
\twocol{0.42}{\vspace{1mm}\FigFlower[50mm]}{0.52}{%
\qline{1} 说出 A 到 F 的名称。
\qline{2} 哪两部分组成 stamen（雄蕊）？
\qline{3} 哪三部分组成 carpel（心皮 / 雌蕊）？
\qline{4} 哪一部分以后变成 fruit（果实）？}
\answers{%
1  A stigma　B style　C ovary\par
　 D anther　E filament　F sepal\par
2  anther + filament（D 和 E）\par
3  stigma + style + ovary（A、B、C）\par
4  ovary（C）；ovule 变成 seed}
\fref{Green Machines p.6}
%%% FRAME 6 {植物繁殖的两种方式}
\renewcommand\arraystretch{1.35}
\begin{tabular}{@{}l l l@{}}
\toprule
 & \textbf{Asexual 无性生殖} & \textbf{Sexual 有性生殖}\\
\midrule
亲本 parents & \gap & 两个\\
gametes（配子）？ & \gap & pollen + ovule\\
后代 offspring & \gap & \gap\\
例子 & runners、bulbs、cuttings & \gap\\
\bottomrule
\end{tabular}\par\vspace{4mm}
（runners 匍匐茎、bulbs 鳞茎、cuttings 扦插）\par\vspace{1mm}
spider plant（吊兰）在长茎上长出 plantlets（小植株）。\par
这是哪种生殖？你怎么知道？
\answers{%
Asexual：one parent｜no gametes｜offspring identical (clones)\par
Sexual：offspring show variation｜例子 flowers and seeds\par
吊兰：asexual，因为 one parent、no gametes、plantlets identical to the parent}
\fref{Green Machines p.3--5}
%%% FRAME 7 {为什么 variation（变异）重要：按顺序写}
\qline{1} 环境改变，比如出现一种新疾病。
\qline{2} 有性生殖产生的后代彼此 \gap[22mm]。
\qline{3} 有些后代有能让它们 \gap[22mm] 的特征。
\qline{4} 这些后代存活并 \gap[22mm]，所以物种 \gap[22mm]。
\par\vspace{3mm}
{\bfseries A}：写第 1--2 句 \quad {\bfseries M}：把 2 和 3 连起来 \quad {\bfseries E}：整条链 + 和无性生殖对比\par\vspace{3mm}
为什么一整片完全一样的植物，会被一种病全部杀死？
\answers{%
2 different　3 survive (the disease)　4 reproduce / continues\par
整片死掉：the plants are clones, so if one gets the disease, they all do（没有一株有抵抗力）}
\fref{Green Machines p.5}
%%% FRAME 8 {从花粉到种子}
把卡片按发生的顺序排好。\par\vspace{2mm}
\qline{A} pollen tube（花粉管）沿着 style（花柱）往下长。
\qline{B} 花粉落到 stigma（柱头）上。
\qline{C} ovule 变成 seed（种子）；ovary（子房）变成果实。
\qline{D} 花粉的细胞核和胚珠的细胞核融合。
\qline{E} 花粉管通过 micropyle（珠孔）进入胚珠。
\qline{F} 昆虫从 anther（花药）带走花粉。
\par\vspace{2mm}
哪几张卡片是 pollination？哪一张是 fertilisation？
\answers{%
顺序：F → B → A → E → D → C\par
pollination = F 和 B　　fertilisation = D}
\fref{Green Machines p.7, p.10}
%%% FRAME 9 {风媒授粉还是虫媒授粉？}
\renewcommand\arraystretch{1.35}
\begin{tabular}{@{}l l l@{}}
\toprule
\textbf{部位} & \textbf{Wind 风媒花} & \textbf{Insect 虫媒花}\\
\midrule
petals 花瓣 & 小、不鲜艳、没香味 & \gap[30mm]\\
stigma 柱头 & \gap[30mm] & 有黏性、在花里面\\
anthers 花药 & 挂在花的外面 & \gap[30mm]\\
pollen 花粉 & \gap[30mm] & 较大、黏或带刺\\
\bottomrule
\end{tabular}\par\vspace{4mm}
满分要写：特征 \textbf{和} 这个特征怎样帮助授粉。\par\vspace{2mm}
为什么草的 anthers 挂在花的外面？
\answers{%
Insect petals：large, brightly coloured, scented, with nectar\par
Wind stigma：large and feathery, hanging outside\par
Insect anthers：firmly attached inside the flower\par
Wind pollen：small, light, made in huge amounts\par
草的 anthers：so the wind can blow the pollen away}
\fref{Green Machines p.8--9}
%%% FRAME 10 {Spot the error (1) 找错误}
草莓用 runners（匍匐茎）和花两种方式产生新植株。比较这两种方式。\par\vspace{2mm}
\samcard{146mm}{\flawed{Runners: the new plant gets gametes from the parent, so it is a bit different from it. Flowers: pollen is carried from the stigma to the anther. This is called fertilisation. The seeds show variation because two parents are involved.}}\par\vspace{2mm}
\textbf{Sam 的答案（英文）有 3 个错误。}找出每一个，改正，并说明为什么会丢分。
\answers{%
1  runners have no gametes: the new plant is identical (a clone)\par
2  pollen goes from the anther to the stigma（Sam 写反了）\par
3  that is pollination, not fertilisation\par
最后一句（两个亲本 → variation）是对的}
\fref{Green Machines p.3--7}
%%% FRAME 11 {Quick review 快速回忆：种子}
\twocol{0.50}{\vspace{2mm}\FigSeed[66mm]}{0.46}{%
\qline{1} 说出 P 到 S 的名称。
\qline{2} 说出 germination（萌发）的三个条件。
\qline{3} 为什么种子要被 dispersed（传播）到离母株很远的地方？}
\answers{%
1  P testa　Q plumule　R radicle　S cotyledon\par
2  water, warmth, oxygen\par
3  so it does not compete with the parent for light, water, minerals and space}
\fref{Green Machines p.11--13}
%%% FRAME 12 {Practice 练习}
\qline{Q1} 百合花有大的白色花瓣、浓香味和 nectar（花蜜）。它是风媒还是虫媒授粉？给两个理由。\slv{A}
\qline{Q2} 种花的人从最好的玫瑰上剪枝 cuttings（扦插），而不是种它的种子。为什么？\slv{M}
\qline{Q3} 水芹种子放进四支试管（20\,\textdegree C）：A 干棉花；B 湿棉花；C 湿棉花放冰箱；D 泡在煮过的水里，上面盖一层油。哪几管会萌发？解释 D 的结果。\slv{E}
\answers{%
Q1  insect：large white petals and scent attract insects; nectar is food for them\par
Q2  cuttings are asexual → identical plants keep the best features; seeds show variation\par
Q3  only B. D has no oxygen, so the seed cannot respire to release energy for growth\par
　　（A 没有水；C 太冷）}
\fref{Green Machines p.5, p.8, p.12}
%%% FRAME 13 {Quick review 快速回忆：叶片结构}
\hspace*{9mm}\FigLeaf[70mm]\hfill
\begin{minipage}[t]{50mm}\vspace{0pt}\raggedright
\qline{1} 说出 A 到 H 的名称。
\qline{2} 哪一层进行最多的 photosynthesis（光合作用）？
\end{minipage}
\answers{%
1  A waxy cuticle　B upper epidermis\par
　 C palisade layer　D spongy layer\par
　 E vein (xylem + phloem)　F lower epidermis\par
　 G guard cell　H stoma\par
2  palisade layer（C）：most chloroplasts}
\fref{Green Machines p.15}
%%% FRAME 14 {Photosynthesis 光合作用：文字方程式}
\vspace{3mm}
\hbox to\textwidth{\hfil\gap[30mm]\ + \gap[22mm]\ \ $\xrightarrow[\rule{20mm}{0.4pt}]{\rule{24mm}{0.4pt}}$\ \ \gap[22mm]\ + \gap[22mm]\hfil}\par\vspace{6mm}
每种原料从哪里进入植物？\par\vspace{2mm}
\qline{} carbon dioxide（二氧化碳）：通过叶子上的 \gap[26mm]
\qline{} water（水）：root hairs（根毛）$\to$ \gap[22mm] $\to$ 叶子\par\vspace{3mm}
chlorophyll（叶绿素）在这个反应里起什么作用？
\answers{%
carbon dioxide + water → glucose + oxygen\par
箭头上：light energy　箭头下：chlorophyll\par
CO$_{2}$ 通过 stomata 进入；water：root hairs → xylem → leaf\par
chlorophyll absorbs light energy}
\fref{Green Machines p.20}
%%% FRAME 15 {Respiration 呼吸作用：同样的物质，反过来}
\vspace{2mm}
\hbox to\textwidth{\hfil glucose + \gap[22mm]\ \ $\longrightarrow$\ \ \gap[30mm]\ + \gap[22mm]\ \ (+ energy 能量)\hfil}\par\vspace{5mm}
\renewcommand\arraystretch{1.35}
\hbox to\textwidth{\hfil\begin{tabular}{@{}l l l@{}}
\toprule
 & \textbf{Photosynthesis 光合作用} & \textbf{Respiration 呼吸作用}\\
\midrule
白天 & \gap[24mm] & \gap[24mm]\\
晚上 & \gap[24mm] & \gap[24mm]\\
\bottomrule
\end{tabular}\hfil}\par\vspace{4mm}
哪个过程在每一个活细胞里、任何时候都在进行？
\answers{%
glucose + oxygen → carbon dioxide + water (+ energy)\par
白天：photosynthesis 进行，respiration 进行\par
晚上：photosynthesis 不进行，respiration 进行\par
答：respiration（每个活细胞，一直都在进行）}
\fref{Green Machines p.23}
%%% FRAME 16 {测试叶子里有没有 starch（淀粉）}
说出每一步的原因。\par\vspace{2mm}
\qline{1} 把叶子放进沸水 1 分钟。
\qline{2} 放进热的 ethanol（乙醇），乙醇管放在一杯热水里。
\qline{3} 用水冲洗。
\qline{4} 滴 iodine（碘液）。变成蓝黑色说明 \gap[34mm]。
\par\vspace{3mm}
为什么乙醇用热水加热，而不是放在 Bunsen（本生灯）火焰上？
\answers{%
1  kills the leaf (stops its reactions)\par
2  removes the chlorophyll, so the colour change can be seen\par
3  softens the leaf / washes off the ethanol\par
4  starch is present\par
ethanol is flammable（易燃），不能靠近火焰}
\fref{Green Machines p.20}
%%% FRAME 17 {Workbook 练习册题：盖锡箔的叶子}
\twocol{0.40}{\vspace{3mm}\FigVarLeaf[56mm]}{0.56}{%
一片 variegated leaf（斑叶，有绿有白）：J 绿色，K 白色。锡箔 L 盖住一部分，在光下放 2 天，然后用碘液测试。\par\vspace{3mm}
哪些部分变蓝黑？逐个解释。}
\answers{%
J（绿色、没盖）：blue-black — chlorophyll + light → photosynthesis → starch\par
K（白色）：stays brown — no chlorophyll\par
L（盖锡箔）：stays brown — no light}
\fref{Green Machines p.22}
%%% FRAME 18 {Spot the error (2) 找错误}
写出光合作用的文字方程式，再说明怎样证明叶子制造了食物。\par\vspace{2mm}
\samcard{146mm}{\flawed{oxygen + water makes glucose + carbon dioxide (light above the arrow). Boil the leaf, put it in hot ethanol, rinse it, then add iodine. If it goes blue-black, the leaf has made glucose. Plants photosynthesise in the day and respire at night.}}\par\vspace{2mm}
\textbf{Sam 的答案（英文）有 3 个错误。}找出每一个，改正，并说明为什么会丢分。
\answers{%
1  carbon dioxide + water → glucose + oxygen（O$_{2}$ 和 CO$_{2}$ 写反了）\par
2  blue-black shows starch, not glucose\par
3  plants respire all the time, day and night\par
实验步骤那一句是对的}
\fref{Green Machines p.20--23}
%%% FRAME 19 {Practice 练习}
\qline{Q1} 两盆植物在光下放 2 天，各罩在一个密封的 bell jar（钟罩）里。1 号罩里还放了一碟 soda lime（碱石灰，会吸收二氧化碳）；2 号罩里放一碟水。各取一片叶子做碘液测试。预测并解释结果。\slv{M}
\qline{Q2} 种满植物的温室里，二氧化碳浓度晚上升高、白天下降。为什么？\slv{E}
\answers{%
Q1  plant 1 stays brown: no carbon dioxide → no photosynthesis → no starch\par
　　plant 2 turns blue-black；jar 2 is the control\par
Q2  night: only respiration, gives out CO$_{2}$\par
　　day: photosynthesis takes in CO$_{2}$ faster than respiration gives it out}
\fref{Green Machines p.21, p.23}
%%% FRAME 20 {Quick review 快速回忆：运输相关的词}
把每个词和它的作用配对。\par\vspace{2mm}
\begin{tabular}{@{}p{50mm}p{96mm}@{}}
\textbf{1} xylem（木质部） & \textbf{A} 把 glucose 从叶子运走\\
\textbf{2} phloem（韧皮部） & \textbf{B} 让气体和水蒸气进出叶子\\
\textbf{3} root hair（根毛） & \textbf{C} 从土壤里吸水\\
\textbf{4} stomata（气孔） & \textbf{D} 把水和 minerals（矿物质）往上运\\
\textbf{5} guard cells（保卫细胞） & \textbf{E} 改变形状，打开和关闭气孔\\
\end{tabular}
\answers{%
1 D　　2 A　　3 C　　4 B　　5 E}
\fref{Green Machines p.16--17}
%%% FRAME 21 {水在植物里怎么走}
\vspace{2mm}
土壤里的水\par
\quad $\downarrow$ 通过 \gap[26mm]\par
根毛细胞\par
\quad $\downarrow$\par
根里和茎里的 \gap[26mm]\par
\quad $\downarrow$\par
叶细胞：水蒸发\par
\quad $\downarrow$\par
以水蒸气的形式从 \gap[26mm] 出去\par\vspace{3mm}
叶子失去水蒸气的过程叫什么？
\answers{%
osmosis　→　xylem　→　stomata\par
叶子失去水蒸气：transpiration（蒸腾作用）}
\fref{Green Machines p.16}
%%% FRAME 22 {水为什么往上走：按顺序写}
芹菜放在蓝色色素水里 2 天，叶子变蓝。\par\vspace{2mm}
\qline{1} 水从叶子通过 stomata \gap[26mm]。
\qline{2} 这会把更多的水沿着 \gap[20mm] 往上 \gap[20mm]。
\qline{3} 色素跟着水走，所以 \gap[20mm] 变蓝。
\qline{4} 把茎切开：蓝点显示 \gap[20mm] 的位置。
\par\vspace{3mm}
{\bfseries A}：色素往上走 \quad {\bfseries M}：说出 xylem \quad {\bfseries E}：transpiration（蒸腾作用）造成拉力
\answers{%
1 evaporates　2 pulls / xylem　3 leaves　4 xylem (vessels)\par
不能写 “the plant sucks / is thirsty”}
\fref{Green Machines p.17}
%%% FRAME 23 {Osmosis 渗透作用}
\twocol{0.32}{\vspace{1mm}\FigMembrane[42mm]}{0.64}{%
左边：糖溶液。右边：水。\par\vspace{2mm}
\qline{1} 哪种形状是水分子？
\qline{2} 水分子往哪边移动？
\qline{3} 为什么糖分子过不去？}
\par\vspace{3mm}
Osmosis 是 \gap[18mm] 从水浓度 \gap[16mm] 的地方到水浓度 \gap[16mm] 的地方，穿过 \gap[30mm] 膜的移动。
\answers{%
1  the small circles（六边形是 sugar）\par
2  right to left, into the sugar solution\par
3  the sugar molecules are too large to fit through the membrane\par
定义空格：water / high / low / semi-permeable}
\fref{Green Machines p.19}
%%% FRAME 24 {Workbook 练习册题：透析袋模型}
\twocol{0.16}{\vspace{1mm}\hspace*{4mm}\FigVisking[9mm]}{0.80}{%
Visking tubing（透析袋）Z 放在一杯水里，30 分钟内液体沿着 X 管上升。\par\vspace{2mm}
\qline{1} 哪一部分相当于根毛？哪一部分相当于茎？
\qline{2} 解释液体为什么上升。
\qline{3} 根毛又长又细，这有什么好处？}
\answers{%
1  Z (Visking tubing) = root hair；X (capillary tube) = stem\par
2  water moves into the sucrose solution by osmosis through the membrane, so the liquid is pushed up\par
3  large surface area, so it absorbs more water}
\fref{Green Machines p.17--18}
%%% FRAME 25 {Spot the error (3) 找错误}
一朵白色康乃馨插在蓝色色素水里，一天后花瓣变蓝。解释原因。\par\vspace{2mm}
\samcard{146mm}{\flawed{The blue water was carried up the phloem tubes. The flower sucks the water up because it is thirsty. Water always gets into a plant by osmosis, from a low concentration of water to a high concentration of water.}}\par\vspace{2mm}
\textbf{Sam 的答案（英文）有 3 个错误。}找出每一个，改正，并说明为什么会丢分。
\answers{%
1  the water goes up the xylem, not the phloem\par
2  not “sucks / thirsty”: water evaporates through the stomata (transpiration) and pulls water up\par
3  osmosis is from a high to a low concentration of water（Sam 写反了）}
\fref{Green Machines p.16--19}
%%% FRAME 26 {Practice 练习}
\qline{Q1} 说出把水送进根毛的过程，以及把水运到叶子的管道。\slv{A}
\qline{Q2} 透析袋模型里，液体一开始在管里 12\,mm 高，10 分钟后 30\,mm，20 分钟后 48\,mm。\par
(a) 计算平均上升速度（每分钟多少 mm）。\par
(b) 解释液体为什么上升。\slv{M}
\qline{Q3} 把植物的根放进很咸的水里，植物会 wilt（萎蔫）。为什么？\slv{E}
\answers{%
Q1  osmosis；xylem\par
Q2 (a)  (48 − 12) ÷ 20 = 1.8 mm per minute（不是 48 ÷ 20）\par
Q2 (b)  water enters the sugar solution by osmosis, so the liquid rises\par
Q3  salty water has a lower concentration of water → water moves out of the roots by osmosis → wilts}
\fref{Green Machines p.16--19}
%%% FRAME 27 {最常见的丢分点：Green Machines}
\qline[11mm]{\Lref{pollfert}} pollination 送花粉；fertilisation 是两个细胞核融合。
\qline[11mm]{\Lref{asexual}} asexual：一个亲本、没有配子、后代一样。
\qline[11mm]{\Lref{adapt}} 写特征，\textbf{也要}写它怎样帮助。
\qline[11mm]{\Lref{photoeq}} light 和 chlorophyll 写在箭头的上下。
\qline[11mm]{\Lref{iodine}} 碘液测的是 starch，不是 glucose。
\qline[11mm]{\Lref{xylem}} xylem 运水；phloem 运 glucose。
\answers{%
让她每条举一个今天出现过的例子：L1 Sam (1)；L2 吊兰；L4、L6 Sam (2)；L7 康乃馨}
\fref{Green Machines p.2}
%%% FRAME 28 {Summary 总结：Green Machines}
\vspace{4mm}
{\renewcommand\baselinestretch{1.3}\sSum
花：pollination 送花粉；fertilisation 合细胞核。\par\vspace{3mm}
种子需要水、温暖、氧气；传播减少竞争。\par\vspace{3mm}
叶子：carbon dioxide + water 用光变成 glucose + oxygen。\par\vspace{3mm}
水：osmosis 进根毛，沿 xylem 上去，transpiration 出去。\par}
\answers{%
本页没有题目。让她闭眼说一遍这四句，再看屏幕对照。}
\fref{Green Machines p.2}
%%% FRAME 29 {Classwork 课堂练习：Green Machines}
在课堂练习纸上做，大约 10 分钟：\par\vspace{3mm}
\qline{A} 花和种子
\qline{B} 光合作用和呼吸作用
\qline{C} 运输和渗透，有一张图要画
\par\vspace{3mm}
每道 \textbf{explain} 题都写完整的英文句子。
\answers{%
做 Classwork Part A–C，约 10 分钟；答案在 classwork\_answers.pdf\par
Part C 的图：0–25 分钟直线上升，之后变平}
\fref{Classwork Part A--C}
%%% FRAME 30 {}
\lessondivider{第二课 \textperiodcentered{} Metallic materials（金属）}{学习目标：把金属的性质和用途联系起来，计算 density（密度），解释 alloys（合金）、rusting（生锈）和金属的化学反应。}
\answers{%
本页没有题目。后半本 booklet 她没写：每页先让她说，再补充。}
\fref{中文版}
%%% FRAME 31 {Do Now 课前小测}
\qline{1} 说出大多数金属共有的两个物理性质。
\qline{2} 说出三种常见的有磁性的金属。
\qline{3} 写出密度的公式。
\qline{4} gallium（镓）在 30\,\textdegree C 熔化。35\,\textdegree C 的天气里它是固体还是液体？
\qline{5} 把金属烧到红热，再慢慢冷却。这种热处理叫什么？
\answers{%
1  任选两个：lustrous, malleable, ductile, sonorous, good conductor, high melting point\par
2  iron, nickel, cobalt\par
3  density = mass ÷ volume\par
4  liquid（35 °C 高于熔点）\par
5  annealing（慢慢冷却）；quenching 是快速冷却}
\fref{Metals p.2--8}
%%% FRAME 32 {The one idea 一个核心思想}
\vspace{6mm}
\hbox to\textwidth{\hfil\FigMetalIdea\hfil}\par\vspace{7mm}
上排：物理性质。下排：化学反应。\par\vspace{2mm}
每道金属题都在这两排中的一排。
\answers{%
问：为什么金子做首饰？答：lustrous（上排）+ very unreactive, does not corrode（下排）}
\fref{Metals p.1}
%%% FRAME 33 {Quick review 快速回忆：性质词}
把每个词和意思配对，然后拼写出来。\par\vspace{2mm}
\begin{tabular}{@{}p{46mm}p{100mm}@{}}
\textbf{1} malleable & \textbf{A} 表面有光泽，能反光\\
\textbf{2} ductile & \textbf{B} 能锤打成形\\
\textbf{3} lustrous & \textbf{C} 敲击时会发出响声\\
\textbf{4} sonorous & \textbf{D} 能拉成细丝\\
\textbf{5} conductor & \textbf{E} 让热或电流通过\\
\end{tabular}
\answers{%
1 B　　2 D　　3 A　　4 C　　5 E\par
拼写：malleable（两个 l，-eable）}
\fref{Metals p.2, p.20}
%%% FRAME 34 {选择金属}
\renewcommand\arraystretch{1.12}
\begin{tabular}{@{}c c c c c@{}}
\toprule
\textbf{金属} & \textbf{熔点} & \textbf{强度} & \textbf{导电性} & \textbf{密度}\\
 & \textbf{(\textdegree C)} & & & \textbf{(g/cm$^{3}$)}\\
\midrule
W & 1450 & 高 & 低 & 7.8\\
X & 650 & 低 & 高 & 2.7\\
Y & 1080 & 中 & 很高 & 8.9\\
Z & 230 & 低 & 中 & 7.3\\
\bottomrule
\end{tabular}\par\vspace{3mm}
哪种金属适合 (a) 桥的框架，(b) 家里的电线，(c) 要求轻的高架电缆？(d) 哪几种会在本生灯火焰（约 800\,\textdegree C）里熔化？
\answers{%
(a) W：high strength, high melting point\par
(b) Y：very high conductivity\par
(c) X：high conductivity and low density (2.7) → light\par
(d) X 和 Z（熔点低于 800 °C）}
\fref{Metals p.6--7}
%%% FRAME 35 {说明选择的理由：按顺序写}
\qline{1} 说出金属的名字。
\qline{2} 说出这个用途需要的性质。
\qline{3} 说明这个性质为什么对这个用途重要。
\qline{4} 比较：为什么不用另一种金属？用数据。
\par\vspace{2mm}
{\bfseries A}：1 + 2 \quad {\bfseries M}：1 到 3 \quad {\bfseries E}：1 到 4，带数字\par\vspace{3mm}
高架电缆我选 \gap[12mm]，因为它 \gap[40mm]。\par
这很重要，因为 \gap[64mm]。\par
Y 导电更好，但是 \gap[64mm]。
\answers{%
I choose X because it has a high electrical conductivity and a low density of 2.7 g/cm$^{3}$.\par
This matters because a light cable puts less weight on the pylons and does not sag.\par
Y conducts better, but its density is 8.9 g/cm$^{3}$, more than three times heavier.}
\fref{Metals p.6--7}
%%% FRAME 36 {Density 密度：公式和体积}
\twocol{0.56}{%
density $= \dfrac{\gap[16mm]}{\gap[16mm]}$\par\vspace{2mm}
单位：g/cm$^{3}$ 或 \gap[16mm]\par\vspace{4mm}
规则形状物体的体积：\par
\quad \gap[50mm]\par\vspace{2mm}
不规则形状物体的体积：\par
\quad 用 \gap[30mm] 法\par\vspace{3mm}
右边物体的体积是多少？}{0.40}{%
\hbox{\FigCylinder{22}{0}\hspace{8mm}\FigCylinder{29}{1}}\par\vspace{1mm}
\hbox to 40mm{\hspace{2mm}放入前\hfil 放入后\hspace{4mm}}}
\answers{%
density = mass ÷ volume；单位 g/cm$^{3}$ 或 g/mL\par
规则物体：length × width × height\par
不规则物体：displacement（排水法）\par
右边物体：29 − 22 = 7 mL}
\fref{Metals p.4}
%%% FRAME 37 {Density 密度：例题}
一把 brass（黄铜）钥匙质量 25.5\,g。放进量筒后，水面从 30.0\,mL 升到 33.0\,mL。\par\vspace{3mm}
\qline[17mm]{第 1 步} 体积 $=$ \gap[16mm] $-$ \gap[16mm] $=$ \gap[16mm] mL
\qline[17mm]{第 2 步} 密度 $=$ \gap[16mm] $\div$ \gap[16mm]
\qline[17mm]{第 3 步} 密度 $=$ \gap[20mm]（写上单位）
\par\vspace{3mm}
检查：答案比 1\,g/mL 大吗？黄铜会下沉，所以应该更大。
\answers{%
第 1 步：33.0 − 30.0 = 3.0 mL\par
第 2 步：25.5 ÷ 3.0\par
第 3 步：8.5 g/mL}
\fref{Metals p.4}
%%% FRAME 38 {Spot the error (4) 找错误}
一个金属螺栓质量 39.5\,g。放进量筒后，水面从 20.0\,mL 升到 25.0\,mL。计算螺栓的密度。\par\vspace{2mm}
\samcard{146mm}{\flawed{volume = 25.0 mL\\density = 39.5 x 25.0 = 987.5 g}}\par\vspace{2mm}
\textbf{Sam 的答案有 3 个错误。}找出每一个，改正，并说明为什么会丢分。
\answers{%
1  volume = 25.0 − 20.0 = 5.0 mL，不是 25.0\par
2  density = mass ÷ volume，不是乘\par
3  unit is g/mL，不是 g\par
正确：39.5 ÷ 5.0 = 7.9 g/mL}
\fref{Metals p.4}
%%% FRAME 39 {Practice 练习}
\qline{Q1} 一块金属 2.0\,cm $\times$ 2.0\,cm $\times$ 3.0\,cm，质量 64.8\,g。计算它的密度。\slv{A}
\qline{Q2} 一枚戒指质量 31.5\,g，放进量筒后水面从 40.0\,mL 升到 43.0\,mL。它是哪种金属？\slv{M}
\par\vspace{1mm}\hspace*{7mm}铁 iron 7.9 \quad 铜 copper 8.9 \quad 银 silver 10.5 \quad 铝 aluminium 2.7（g/mL）
\qline{Q3} 银的导电性比铜好。解释为什么家里的电线用铜做。\slv{E}
\answers{%
Q1  2.0 × 2.0 × 3.0 = 12.0 cm$^{3}$；64.8 ÷ 12.0 = 5.4 g/cm$^{3}$\par
Q2  43.0 − 40.0 = 3.0 mL；31.5 ÷ 3.0 = 10.5 g/mL → silver\par
Q3  copper is almost as good a conductor, much cheaper, and ductile (can be drawn into wires)}
\fref{Metals p.4, p.6--7}
%%% FRAME 40 {Quick review 快速回忆：热处理}
\qline{1} 烧到红热，然后放进冷水里快速冷却。
\qline{2} 烧到红热，然后慢慢冷却。
\qline{3} 强烈加热，直到出现蓝色光泽。
\par\vspace{2mm}
说出名称：annealing（退火）、quenching（淬火）或 tempering（回火）。\par
再配上结果：\textbf{A} 软、容易弯 \quad \textbf{B} 硬但脆 \quad \textbf{C} 有韧性、有弹性\par\vspace{3mm}
\qline{4} 发夹要能弹回原来的形状，需要哪种热处理？
\answers{%
1  quenching → B（硬但脆）\par
2  annealing → A（软）\par
3  tempering → C（有韧性、有弹性）\par
4  tempering（发夹要 springy）}
\fref{Metals p.8}
%%% FRAME 41 {合金为什么更硬 (1)}
纯金属：每个原子都一样大，排成整齐的层。\par\vspace{5mm}
\hbox to\textwidth{\hfil\FigPure[5mm]\hspace{14mm}\FigPureSlid[5mm]\hfil}\par\vspace{2mm}
\hbox to\textwidth{\hfil 推之前\hspace{46mm}推一下以后\hspace{12mm}\hfil}\par\vspace{4mm}
推的时候，这些层怎么了？所以纯金属是硬还是软？
\answers{%
the layers slide over each other easily\par
so a pure metal is soft（所以 malleable、ductile）}
\fref{Metals p.2, p.9}
%%% FRAME 42 {合金为什么更硬 (2)}
alloy（合金）：一种金属和另一种元素混合。\par\vspace{1mm}
\twocol{0.40}{\vspace{1mm}\FigAlloy[4.4mm]}{0.56}{%
\qline{1} 原子的大小 \gap[28mm]。
\qline{2} 这 \gap[28mm] 了整齐的层。
\qline{3} 层不能轻易地 \gap[28mm]。
\qline{4} 所以合金比纯金属 \gap[22mm]。}
\par\vspace{3mm}
{\bfseries A}：更硬 \quad {\bfseries M}：+ 原子大小不同 \quad {\bfseries E}：整条链
\answers{%
1 different　2 disrupts　3 slide over each other　4 harder}
\fref{Metals p.9}
%%% FRAME 43 {合金的用途}
\renewcommand\arraystretch{1.3}
\begin{tabular}{@{}l l l@{}}
\toprule
\textbf{合金} & \textbf{主要金属} & \textbf{用途}\\
\midrule
bronze 青铜 & 铜 + \gap[18mm] & 齿轮、机器轴承\\
solder 焊锡 & 锡 + 铅 & \gap[44mm]\\
stainless steel 不锈钢 & 铁、铬、镍 & \gap[44mm]\\
nitinol 镍钛合金 & 镍 + 钛 & \gap[44mm]\\
\bottomrule
\end{tabular}\par\vspace{4mm}
为什么焊锡适合连接电线？\par
合金是 compound（化合物）还是 mixture（混合物）？
\answers{%
bronze：tin　　solder：joining electrical wires\par
stainless steel：cutlery, sinks（不生锈）　　nitinol：braces, glasses frames（shape memory）\par
solder：low melting point　　alloy 是 mixture（混合物）}
\fref{Metals p.9, p.11}
%%% FRAME 44 {合金的熔点}
\twocol{0.50}{\vspace{1mm}\hspace*{10mm}\FigMPGraph[4.6mm]}{0.46}{%
两种金属 X 和 Y 的混合物。\par\vspace{2mm}
\qline{1} 纯 X 的熔点是多少？纯 Y 呢？
\qline{2} 哪个混合比例的熔点最低？
\qline{3} 描述从 0\,\% 到 70\,\% X 的变化趋势。}
\answers{%
1  pure X 262 °C；pure Y 330 °C\par
2  70 \% X，about 190 °C\par
3  the melting point falls steadily as the percentage of X increases}
\fref{Metals p.10--11}
%%% FRAME 45 {Practice 练习}
\qline{Q1} 哪种热处理让钢变硬但变脆？\slv{A}
\qline{Q2} brass（黄铜）是铜和锌的合金。解释为什么黄铜比纯铜硬。\slv{M}
\qline{Q3} 汽车的弹簧和车门板都用钢做。解释它们各应该做哪种热处理。\slv{E}
\answers{%
Q1  quenching\par
Q2  zinc atoms are a different size → disrupt the layers → layers cannot slide easily\par
Q3  spring: tempering (tough, springy)；door panel: annealing (soft, easy to press into shape)}
\fref{Metals p.8--9}
%%% FRAME 46 {Quick review 快速回忆：腐蚀和生锈}
判断对错（true / false）：\par\vspace{2mm}
\qline{1} rusting（生锈）是 corrosion（腐蚀）的一种。
\qline{2} 任何金属都会生锈。
\qline{3} 盐水让铁生锈得更快。
\qline{4} 金在空气里很快被腐蚀。
\qline{5} 腐蚀是金属和氧气、水等物质发生的化学反应。
\answers{%
1 true　2 false（只有 iron/steel 会 rust）　3 true　4 false　5 true}
\fref{Metals p.13, p.21}
%%% FRAME 47 {生锈实验 (1)：装置}
\twocol{0.54}{\vspace{2mm}\FigRustTubes[2.8mm]}{0.43}{%
铁钉放 3 天。calcium chloride = 氯化钙。\par\vspace{2mm}
\qline{A} 氯化钙吸走了 \gap[20mm]。
\qline{B} 煮沸赶走了水里的 \gap[20mm]；油阻止 \gap[20mm]。
\qline{C} 钉子既有 \gap[14mm] 又有 \gap[14mm]。}
\answers{%
A  water (moisture)\par
B  oxygen；air getting back in\par
C  water and oxygen}
\fref{Metals p.17}
%%% FRAME 48 {生锈实验 (2)：按顺序写}
3 天后的结果： \quad A \gap[22mm] \quad B \gap[22mm] \quad C \gap[22mm]\par\vspace{4mm}
\qline{1} 只有 \gap[10mm] 号试管生锈。
\qline{2} A 管有空气，但没有 \gap[22mm]。
\qline{3} B 管有水，但没有 \gap[22mm]。
\qline{4} 所以铁生锈需要 \gap[50mm]。
\par\vspace{3mm}
{\bfseries A}：C 生锈 \quad {\bfseries M}：+ A 和 B 各缺什么 \quad {\bfseries E}：用三支试管得出结论
\answers{%
结果：A no rust　B no rust　C rust\par
1 C　2 water　3 oxygen　4 both water and oxygen}
\fref{Metals p.17}
%%% FRAME 49 {Spot the error (5) 找错误}
钢制回形针分别封在三支试管里：P 有氯化钙（干燥空气），Q 在煮过的水里、上面有一层油，R 在自来水里、上面有空气。描述并解释结果。\par\vspace{2mm}
\samcard{146mm}{\flawed{P: no rust, because there is no air. Q: a little rust, because it has water. R: lots of rust. So iron only needs water to rust.}}\par\vspace{2mm}
\textbf{Sam 的答案（英文）有 3 个错误。}找出每一个，改正，并说明为什么会丢分。
\answers{%
1  P has air; it does not rust because it has no water\par
2  Q does not rust: boiled water has no oxygen and the oil keeps air out\par
3  iron needs both water and oxygen\par
R 那一句是对的}
\fref{Metals p.13, p.17}
%%% FRAME 50 {防止生锈}
\renewcommand\arraystretch{1.35}
\begin{tabular}{@{}l l l@{}}
\toprule
\textbf{物体} & \textbf{方法} & \textbf{为什么用这种方法}\\
\midrule
钢栅栏 & \gap[26mm] & 面积大、便宜、容易重做\\
门铰链 & 上油或润滑脂 & \gap[40mm]\\
垃圾桶 & \gap[26mm] & 镀一层锌\\
厨房水槽 & \gap[26mm] & 永远不需要涂层\\
\bottomrule
\end{tabular}\par\vspace{4mm}
每种方法都是让 \gap[18mm] 和 \gap[18mm] 接触不到铁。
\answers{%
钢栅栏：painting\par
门铰链：moving parts, paint would rub off\par
垃圾桶：galvanising\par
厨房水槽：stainless steel\par
最后一句：water and oxygen}
\fref{Metals p.13, p.15}
%%% FRAME 51 {Practice 练习}
\qline{Q1} 铁生锈需要哪两种物质？\slv{A}
\qline{Q2} 为什么海边小镇的汽车，比内陆同样的汽车生锈更快？\slv{M}
\qline{Q3} 设计一个实验：盐水是否比自来水让铁钉生锈更快？说出你改变什么、测量什么、保持什么不变。\slv{E}
\answers{%
Q1  water and oxygen\par
Q2  sea air carries salt, and salt speeds up rusting\par
Q3  change: tap water or salt water；measure: amount of rust (or mass) after 3 days\par
　　keep the same: nail type and size, volume of water, temperature, time；repeat}
\fref{Metals p.13, p.17}
%%% FRAME 52 {Quick review 快速回忆：文字方程式的规律}
\qline{1} 金属 + oxygen（氧气）$\to$ \gap[40mm]
\qline{2} 金属 + water（水）$\to$ \gap[34mm] + \gap[26mm]
\qline{3} 金属 + hydrochloric acid（盐酸）$\to$ \gap[34mm] + \gap[26mm]
\qline{4} 检验 hydrogen（氢气）：点燃的木条发出 \gap[30mm]。
\answers{%
1  metal oxide\par
2  metal hydroxide + hydrogen\par
3  metal chloride + hydrogen\par
4  squeaky pop}
\fref{Metals p.12, p.18--19}
%%% FRAME 53 {写文字方程式}
\qline{(a)} 钢丝绒在氧气里燃烧。\par\hspace*{7mm}iron + oxygen $\to$ \gap[40mm]
\qline{(b)} lithium（锂）在冷水里冒泡。\par\hspace*{7mm}lithium + water $\to$ \gap[40mm] + \gap[22mm]
\qline{(c)} 铝粉放进稀盐酸。\par\hspace*{7mm}\gap[120mm]
\qline{(d)} 金放进稀盐酸。\par\hspace*{7mm}\gap[60mm]
\answers{%
(a) iron oxide\par
(b) lithium hydroxide + hydrogen\par
(c) aluminium + hydrochloric acid → aluminium chloride + hydrogen\par
(d) no reaction}
\fref{Metals p.12, p.18--19}
%%% FRAME 54 {The reactivity series 金属活动性顺序}
\vspace{3mm}
\hbox to\textwidth{\hfil\FigReactivity\hfil}\par\vspace{4mm}
\qline{1} tin（锡）会和稀酸反应吗？
\qline{2} 为什么地下能找到纯的金，却找不到纯的铁？
\answers{%
1  yes（tin 在 react with acids 的范围里，很慢）\par
2  gold is very unreactive, so it stays as the metal；iron is more reactive, so it is found as compounds (ores)}
\fref{Metals p.20}
%%% FRAME 55 {Spot the error (6) 找错误}
完成每个文字方程式，或者写 ``no reaction''（不反应）。\par\vspace{2mm}
\samcard{146mm}{\flawed{1. calcium + oxygen makes calcium oxygen\\
2. iron + hydrochloric acid makes iron chloride + water\\
3. potassium + water makes potassium hydroxide + hydrogen\\
4. silver + hydrochloric acid makes silver chloride + hydrogen}}\par\vspace{2mm}
\textbf{Sam 的答案（英文）有 3 个错误。}找出每一个，改正，并说明为什么会丢分。
\answers{%
1  calcium oxide，不是 calcium oxygen\par
2  iron chloride + hydrogen，不是 water\par
3  silver + hydrochloric acid → no reaction\par
第 3 条（potassium）是对的}
\fref{Metals p.12, p.18--20}
%%% FRAME 56 {Practice 练习}
\qline{Q1} 写出铝和氧气反应的文字方程式。\slv{A}
\qline{Q2} 四种金属放进稀酸。按活泼程度从高到低排列，并给出证据。\slv{M}
\par\vspace{1mm}\hspace*{7mm}\begin{tabular}{@{}l l@{}}
P & 气泡很多、很快\\
Q & 没有气泡\\
R & 气泡很少、很慢\\
S & 连冷水都能反应\\
\end{tabular}
\qline{Q3} 描述一个用稀盐酸证明锌比铜活泼的实验，并说出你预期的结果。\slv{E}
\answers{%
Q1  aluminium + oxygen → aluminium oxide\par
Q2  S, P, R, Q — S reacts even with cold water; P many bubbles fast; R few, slow; Q none\par
Q3  same size zinc and copper, same volume and concentration of acid：zinc fizzes (hydrogen), copper no reaction}
\fref{Metals p.18--20}
%%% FRAME 57 {最常见的丢分点：Metals}
\qline[11mm]{\Lref{density}} 密度 = 质量 $\div$ 体积，单位 g/mL 或 g/cm$^{3}$。
\qline[11mm]{\Lref{displace}} 体积 = 放入后的读数 $-$ 放入前的读数。
\qline[11mm]{\Lref{justify}} 性质 + 为什么重要 + 对比。
\qline[11mm]{\Lref{alloy}} 原子大小不同，层不能滑动。
\qline[11mm]{\Lref{rust}} 生锈需要水 \textbf{和} 氧气。
\qline[11mm]{\Lref{wordeq}} oxide；hydroxide + hydrogen；chloride + hydrogen。
\answers{%
让她每条举例：L10、L11 Sam 的螺栓；L13 黄铜；L14 回形针；L15 Sam 的方程式}
\fref{Metals p.1}
%%% FRAME 58 {Summary 总结：Metallic materials}
\vspace{4mm}
{\renewcommand\baselinestretch{1.3}\sSum
性质决定用途：说出性质、为什么重要，再比较。\par\vspace{3mm}
密度 = 质量 $\div$ 体积；用排水法测体积。\par\vspace{3mm}
合金更硬：大小不同的原子让层不能滑动。\par\vspace{3mm}
生锈需要水和氧气；活泼程度决定金属怎样反应。\par}
\answers{%
本页没有题目。先闭眼说一遍，再看屏幕。}
\fref{Metals p.20}
%%% FRAME 59 {Classwork 课堂练习：Metallic materials}
在课堂练习纸上做，大约 10 分钟：\par\vspace{3mm}
\qline{D} 性质和密度
\qline{E} 合金和热处理
\qline{F} 腐蚀和反应，有一个实验要设计
\par\vspace{3mm}
计算题每一步都写出来，带单位。
\answers{%
做 Classwork Part D–F，约 10 分钟；答案在 classwork\_answers.pdf\par
Part D 密度题：2.7 g/cm$^{3}$（aluminium）}
\fref{Classwork Part D--F}
%%% FRAME 60 {Mini mock 小模考}
\vspace{4mm}
{\sEmph\bfseries 17 分钟 \textperiodcentered{} 两个单元 \textperiodcentered{} 写在卷子上}\par\vspace{6mm}
每一问都要回答。\par\vspace{2mm}
每题按 A、M、E 评分。\par\vspace{2mm}
Explain 题：按今天练习的顺序写。
\answers{%
发 mock.pdf，计时 17 分钟，不提示；答案在 mock\_answers.pdf}
\fref{Mini mock}
%%% FRAME 61 {Homework 作业}
\qline[12mm]{hw1} Green Machines，大约 40 分钟
\qline[12mm]{hw2} Metallic materials，大约 40 分钟
\par\vspace{3mm}
另外：Metals booklet 第 10--11 页，画焊锡的图，并回答 Q1--2。\par\vspace{3mm}
在考试前，分两天做完。
\answers{%
hw1、hw2 各约 40 分钟；答案在 hw1\_answers.pdf、hw2\_answers.pdf}
\fref{hw1, hw2}
